\section{Appendici}
\begin{frame}{Perché il quantum computing}
	\begin{minipage}{\linewidth}
		\begin{minipage}{.3\linewidth}
			\begin{block}{Grazie a}
				\begin{itemize}
					\item Superposition
					\item Entanglement
				\end{itemize}
			\end{block}
		\end{minipage}
		\hfill
		\begin{minipage}{.68\linewidth}
			\begin{gather*}
				\mathcal{H}(t) \Ket{\psi(t)} = i \hbar \frac{\partial}{\partial t} \Ket{\psi(t)}\\ 
				\Braket{A_0 \otimes B_0} + \Braket{A_0 \otimes B_1} + \Braket{A_1 \otimes B_0} - \Braket{A_1 \otimes B_1} = 2\sqrt{2}
			\end{gather*}
		\end{minipage}
	\end{minipage}
	\begin{minipage}{\linewidth}
		\begin{minipage}{.3\linewidth}
			\begin{block}{Violiamo}
				\begin{center}
					Strong Church-Turing Thesis
				\end{center}
			\end{block}
		\end{minipage}
		\hfill
		\begin{minipage}{.68\linewidth}
			\vspace{.5cm}
			\begin{figure}
				\centering
				\includegraphics[width=.6\linewidth]{ntm}
				\caption{Probabilistic Turing Machine}
			\end{figure}
		\end{minipage}
	\end{minipage}
\end{frame}

\begin{frame}{Quantum Gate}
	\begin{minipage}{.35\linewidth}
		\begin{block}{}
			\begin{itemize}
				\item Formalismo molto studiato
				\item Gate Reversibili
				\item  Set di porte universali
				\item  Algoritmi di Shore, ecc.
			\end{itemize}
		\end{block}
	\end{minipage}
	\hfill
	\begin{minipage}{.6\linewidth}
		\begin{figure}
			\includegraphics[width=\linewidth]{Shor}
			\caption{Step for Shor algoritm (Thomas G. Wong)}
		\end{figure}
	\end{minipage}
\end{frame}


\begin{frame}{Quantum Annealer}
	\begin{minipage}{.48\linewidth}
		\begin{figure}
			\includegraphics[width=.65\linewidth]{QA_fun}
			\caption{Complete sovrapposition $\rightarrow$ ground state $\mathcal{H}_f$}
		\end{figure}
	\end{minipage}
	\hfill
	\begin{minipage}{.48\linewidth}
		\begin{block}{}
			\begin{itemize}
				\item Problemi in forma QUBO
				\item Ispirato a Simulated Annealling
				\item Effetto tunneling
				\item Adiabatic Quantum Computing (AQC)
				\item $\mathcal{H}(t) =s(t) \mathcal{H}_i + (i-s(t)) \mathcal{H}_f $
			\end{itemize}
		\end{block}
	\end{minipage}	
\end{frame}

\begin{frame}[fragile]{QMASM}{Esempio: Macro}
	\begin{minipage}{.3\linewidth}
		\begin{figure}[h]
			\begin{minipage}{\linewidth}
				\hrule
				\begin{minted}{text}
					
# Y = A OR B
!begin_macro OR
	$A  0.5
	$B  0.5
	$Y -1
	
	$A $B  0.5
	$A $Y -1
	$B $Y -1
!end_macro OR
				\end{minted}
				\hrule
			\end{minipage}
			\caption{\textbf{or} gate}
		\end{figure}
	\end{minipage}
	\hfill
	\begin{minipage}{.3\linewidth}
		\begin{figure}[h]
			\begin{minipage}{\linewidth}
				\hrule
				\begin{minted}{text}
					
# Y = NOT A
	!begin_macro NOT
	$A $Y 1.0
!end_macro NOT
				\end{minted}
				\hrule
			\end{minipage}
			\caption{\textbf{not} gate}
		\end{figure}
	\end{minipage}
	\hfill
	\begin{minipage}{.3\linewidth}
		\begin{figure}[h]
			\begin{minipage}{\linewidth}
				\hrule
				\begin{minted}{text}

# Y = A AND B
!begin_macro AND
	$A -0.5
	$B -0.5
	$Y  1
	
	$A $B  0.5
	$A $Y -1
	$B $Y -1
!end_macro AND
				\end{minted}
				\hrule
			\end{minipage}
			\caption{\textbf{and} gate}
		\end{figure}
	\end{minipage}
	\begin{block}{}
		\centering
		possiamo racchiudere queste macro nel file \texttt{gates.qmasm}
	\end{block}
\end{frame}

\begin{frame}[fragile]{QMASM}{Esempio: $y = x_1 \land \neg(x_2 \lor x_3)$}
	\begin{center}
		\begin{minipage}{.25\linewidth}
			\begin{figure}[h]
				\begin{minipage}{\linewidth}
					\hrule
					\begin{minted}{text}
						
!include <gates>

!use_macro OR x2_or_x3
	x2_or_x3.$A = x2
	x2_or_x3.$B = x3
	x2_or_x3.$Y = $x4

!use_macro NOT not_x4
	not_x4.$A = $x4
	not_x4.$Y = $x5

!use_macro AND x1_and_x5
	x1_and_x5.$A = x1
	x1_and_x5.$B = $x5
	x1_and_x5.$Y = y
					\end{minted}
					\hrule
				\end{minipage}
				\caption{CircSat problem}
			\end{figure}
		\end{minipage}
		\hspace{1cm}
		\begin{minipage}{.6\linewidth}
			\begin{block}{}
				\say{Pinniamo} il valore di $y$ per ottenere l'assegnamento delle $x_i$ che verificano la formula logica:
				\begin{center}
					\scriptsize \texttt{qmasm --run --pin="y := true" circsat.qmasm}
				\end{center}
			\end{block}
			\begin{figure}[h]
				\centering
				\begin{minipage}{.72\linewidth}
					\hrule
					\begin{minted}{text}
						
Solution #1 (energy = -20.0000, tally = 647):

							Variable  Value
							--------  -----
							x1        True
							x2        False
							x3        False
							y         True
					\end{minted}
					\hrule
				\end{minipage}
				\caption{CircSat solution}
			\end{figure}
		\end{minipage}
	\end{center}
\end{frame}

\begin{frame}[fragile]{Yosys - edif2qmasm}{Esempio: moltiplicazione tra interi}
	\begin{minipage}{.48\linewidth}
		\begin{figure}[h]
			\begin{minipage}{\linewidth}
				\hrule
				\begin{minted}{Verilog}
					
module mult (multiplicand, multiplier, product);
	input [1:0] multiplicand;
	input [1:0] multiplier;
	output[2:0] product;
	
	assign product = multiplicand * multiplier;
endmodule
				\end{minted}
				\hrule
			\end{minipage}
			\caption{Factorization problem}
		\end{figure}
		\vspace{-.7cm}
		\begin{block}{Traduzione in EDIF}
			\begin{center}
				\scriptsize \texttt{yosys myfile.v synth.ys -b edif -o myfile.edif}
			\end{center}
		\end{block}
		\begin{block}{Traduzione in QMASM}
			\begin{center}
				\scriptsize \texttt{edif2qmasm -o="myfile.qmasm" myfile.edif}
			\end{center}
		\end{block}
	\end{minipage}
	\hfill
	\begin{minipage}{.48\linewidth}
		\begin{block}{Esecuzione}
			\begin{center}
				\scriptsize \texttt{qmasm --run --pin="mult.product[2:0] := 110" --solver="sim\_anneal" mult.qmasm}
			\end{center}
		\end{block}
		\vspace{-.1cm}
		\begin{figure}[h]
			\begin{minipage}{.9\linewidth}
				\hrule
				\begin{minted}{text}
					
Solution #1 (energy = -57.5000, tally = 68):

				Variable              Value
				--------------------  -----
				mult.multiplicand[0]  False
				mult.multiplicand[1]  True
				mult.multiplier[0]    True
				mult.multiplier[1]    True
				mult.product[0]       False
				mult.product[1]       True
				mult.product[2]       True
				\end{minted}
				\hrule
			\end{minipage}
			\caption{Factorization Solution}
		\end{figure}
	\end{minipage}
\end{frame}

\begin{frame}[fragile]{Update al progetto}{Incompatibilità tra output edif2qmasm e input QMASM}
	\begin{minipage}{.31\linewidth}
		\begin{figure}
			\begin{minipage}{.95\linewidth}
				\hrule
				\begin{minted}{verilog}
					
// Define hates(atom, atom).
module \hates/2 (A, B, Valid);
	...
endmodule
				\end{minted}
				\hrule
			\end{minipage}
			\caption{Verilog}
		\end{figure}
		\vspace{-.6cm}
		\begin{figure}
			\begin{minipage}{.95\linewidth}
				\hrule
				\begin{minted}{text}
					
# hates/2
!begin_macro id00011
	...
!end_macro id00011

# enemies/2
!begin_macro id00010
!use_macro OR $id00031
!use_macro hates/2
...
				\end{minted}
				\hrule
				\caption{qmasm}
			\end{minipage}
		\end{figure}
	\end{minipage}
	\hfill
	\begin{minipage}{.67\linewidth}
		\begin{figure}
			\begin{minipage}{.95\linewidth}
				\hrule
				\begin{minted}{python}
					
with open(file, 'r') as input:
first = input.readline()
second = input.readline()
while(second_row != ""):
	if first.startswith("#") and second.startswith("!begin_macro"):
		self.name[first[2:-1]] = second[len("!begin_macro "):-1]
	first_row = second_row
	second_row = input.readline()

with open(file, 'r') as input:
doc = input.read()
lines = doc.splitlines()
for line in lines:
	for word in line.split():
		if word in self.name.keys():
			line = line.replace(word, self.name[word])
	self.new_lines.append(line)
				\end{minted}
				\hrule
			\end{minipage}
			\caption{Preprocessing}
		\end{figure}
	\end{minipage}
\end{frame}

\begin{frame}{Ontologie}
	\begin{minipage}{.48\linewidth}
		\begin{block}{Cos'è un'ontologia}
			\begin{itemize}
				\item Insieme di fatti
				\item Riguardanti un dominio di interesse
				\item Prevengono interpretazioni sbagliate
				\item Assicurano la cooperazione tra software
			\end{itemize}
		\end{block}
		\begin{block}{Come si definisce un'ontologia}
			\begin{itemize}
				\item Ontology Web Language (OWL)
				\item Diversi \say{flavours}
				\begin{itemize}
					\item OWL Full
					\item OWL DL
				\end{itemize}
				\item Classi Individui e Relazioni
			\end{itemize}
		\end{block}
	\end{minipage}
	\hfill
	\begin{minipage}{.48\linewidth}
		\begin{figure}[h]
			\includegraphics[width=\linewidth]{europeana_relation.pdf}
		\end{figure}
	\end{minipage}
\end{frame}

\begin{frame}{Inferenze in OWL}
	\begin{minipage}{.48\linewidth}
		\begin{block}{Complessità}
			\begin{itemize}
				\item Variazione sintattica di $\mathcal{SROIQ}$
				\item $\mathcal{ALC}$ è una restrizione di $\mathcal{SROIQ}$
				\item $\mathcal{ALC}$ è PSpace-hard
			\end{itemize}
		\end{block}
		\begin{block}{Reasoner}
			\begin{itemize}
				\item Pellet
				\item Fact++
				\item HermiT
				\item \dots
			\end{itemize}
		\end{block}
	\end{minipage}
	\hfill
	\begin{minipage}{.48\linewidth}
		\begin{figure}[h]
			\includegraphics[width=.8\linewidth]{Complexity.pdf}
		\end{figure}
	\end{minipage}
\end{frame}